%% kapa: exam=1 round=24 year=2013 period=1 form=A subject=civil_law subject_ko=민법 seq=12 number=12
%% source: Q-Net 감정평가사 기출문제 (cst00309) artlSeq=5073802 file=extracted/1008/5073802/2013년 제24회 감정평가사 1차시험 기출문제/1교시(A형).hwp
%% conversion: kordoc 4.24.51 -> tools/gen1.py (auto); review=proofread-sonnet; flags=none
\begin{question}{12}
착오에 관한 설명으로 옳지 않은 것은? (다툼이 있으면 판례에 의함)
\begin{choices}
  \item 착오가 표의자의 중대한 과실로 인한 때에는 착오를 이유로 취소할 수 없다.
  \item 소취하 등 소송행위에도 민법상 사기 또는 착오로 인한 취소규정이 적용된다.
  \item 계약당사자 쌍방이 계약의 기초가 되는 사항에 관하여 같은 내용으로 착오하여 그 사항에 관하여 약정하지 않은 경우 그 착오가 없었다면 약정하였을 내용으로 당사자의 의사를 보충할 수 있다.
  \item 착오를 이유로 취소한 자는 이로 인하여 상대방에게 손해가 발생했더라도 불법행위책임을 부담하지 않는다.
  \item 장래 도시계획이 변경되어 호텔의 신축허가를 받을 수 있을 것이라고 스스로 생각하여 토지의 매매계약을 체결했으나 그 후 예상대로 되지 않은 경우 착오를 이유로 취소할 수 없다.
\end{choices}
\end{question}
\answer{2}
